\section{The real cyclotomic subfield and the empty norm range}
\label{sec:literal-small}

We first prove the unconditional conclusion, stating the precise source
data that enter it. Write $F$ for the field called the host field in
\citet{AlpogeLean2026}. A frame gives a finite abelian group $A$ of
order $n$, a surjection
\[
 h:A\longrightarrow H_e=(\Z/2^{e+2}\Z)^\times/\{\pm1\},
 \qquad a_0\in A,
                                                        \tag{2.1}
\]
and the implication $h(ja_0)=0\Rightarrow ja_0=0$ for every integer
$j$. Put $n_0=\operatorname{ord}(a_0)$ and $o=n/n_0$.
The full host has $e\ge2$, $o^{64}\le2^e$, and
$\Gal(F/\Q)\simeq C_2^2\times A$.
These are the source declarations \texttt{Frame}, \texttt{Ranks.n},
and \texttt{Host.o\_pow}; the order $n$ is the same parameter as
in its degree window.

There is more information than an abstract cyclic quotient. The
host includes an embedding
\[
 \iota:F\hookrightarrow C=\Q(\zeta_N),\qquad
 N=Q_Sqd,\qquad 2^{e+2}\mid Q_S,
                                                        \tag{2.2}
\]
and an Artin map from $(\Z/N\Z)^\times$ onto $\Gal(F/\Q)$.
The source equality \texttt{art\_spec} says that this map is the
restriction of the cyclotomic action through $\iota$.
The equality \texttt{art\_h} identifies its projection to $H_e$
with reduction modulo $2^{e+2}$ followed by quotienting by $\{\pm1\}$.
We identify $F$ with its image under $\iota$.

\begin{lemma}[A common subfield]\label{lem:actual-cyclotomic}
Every such host contains
\[
       B_e=\Q(\zeta_{2^{e+2}})^+,
       \qquad B=\Q(\zeta_{16})^+\subset B_e.
                                                        \tag{2.3}
\]
\end{lemma}
\begin{proof}
An automorphism of $C$ is given by a unit $u$ modulo $N$.
If it fixes $F$, the restriction equality in (2.2) and injectivity
of $\iota$ make its Artin image the identity. Its projection to
$H_e$ is therefore the identity. By \texttt{art\_h}, this means
$u\equiv\pm1\pmod{2^{e+2}}$, exactly the condition that the
automorphism fix $B_e$. Hence
$\Gal(C/F)\subset\Gal(C/B_e)$, and the Galois correspondence
gives $B_e\subset F$. The second inclusion follows from $e\ge2$.
\end{proof}

\begin{lemma}[Degree forces a large cyclotomic conductor]
\label{lem:large-tower}
For the full host,
\[
 n\le2^{65e/64},\qquad
 M:=2^{e+2}\ge4n^{64/65}.
                                                        \tag{2.4}
\]
If $x>1$, $\epsilon>0$, and
$n\ge x^{2+8\epsilon-\epsilon/100}$, then $M>2x+1$.
\end{lemma}
\begin{proof}
The implication following (2.1) makes $h$ injective on
$\langle a_0\rangle$: apply it to the difference of any two
elements with equal image. The group $H_e$ has order $2^e$, so
$n_0\le2^e$. Also $n_0\mid n$ by Lagrange's theorem and thus
$n=o n_0$. Combining this with $o^{64}\le2^e$ proves (2.4).
For completeness, $|H_e|=2^e$ because the group of units modulo
$2^{e+2}$ has order $2^{e+1}$ and $-1$ has order two.
The lower degree window now gives
\[
 M\ge4x^{(64/65)(2+799\epsilon/100)}>4x>2x+1.
                                                        \tag{2.5}
\]
The exponent exceeds one, and the final inequality uses $x>1$.
\end{proof}

\begin{proof}[Proof of Theorem~\ref{thm:empty-norms}]
Let $\mathfrak P$ be a prime of any extension $F'/F$ above an odd
rational prime $p$. Write $\N\mathfrak P=p^f$. Its restriction
to $B_e$ has residue degree dividing $f$, since the corresponding
residue field embeds in that at $\mathfrak P$. The field $B_e$
is unramified at $p$; its Frobenius is the class of $p$ modulo
$M$, up to sign. Consequently
\[
                         p^f\equiv\pm1\pmod M.
                                                        \tag{2.6}
\]
If $p^f\le2x$, (2.5) gives $1<p^f<M-1$, contradicting (2.6).
This argument does not assume that $F$ or $F'$ is unramified at $p$.

The logarithmic derivative of a Hecke Euler product is a sum over
powers of prime-ideal norms. A coefficient of index at most $2x$
and odd residue characteristic would therefore require a prime
ideal of norm at most $2x$, which we have excluded. Deleting the
factors over $2$ proves the coefficientwise vanishing assertion.
All its finite linear combinations vanish as well. At $p=7$ the
cubic trace is $-1$, because $269\equiv3\pmod7$ is not a cube.
Thus once $x\ge7/2$ these coefficients cannot recover all the
target coefficients through $2x$.
\end{proof}

The exclusion is stronger than the fixed common-subfield obstruction.
Even without the degree window, (2.3) implies that every odd prime
$p\not\equiv\pm1\pmod{16}$ has no degree-one prime in a literal
host. Among split-$K$ classes modulo $48$, the only potentially
visible classes are $1,31$; the missing classes are
$7,13,19,25,37,43$. With the full degree window every odd prime
in the range is missing, including the two formerly visible classes.

\subsection{Why this is not an obstruction to the source's theta gain}

There is also a positive consequence for the local orbit mechanism.
For an odd prime $p\le2x$, let
\[
 v=\begin{cases}v_2(p-1),&p\equiv1\pmod4,\\
                 v_2(p+1),&p\equiv3\pmod4.
   \end{cases}
\]
Choose $u=p$ or $u=-p$ so that $u\equiv1\pmod4$. Then
$v_2(u-1)=v\ge2$, and successive factorizations of
$u^{2^{j+1}}-1$ show
$v_2(u^{2^j}-1)=v+j$: the second factor always has valuation one.
A power of $u$ is $1$ modulo four, so it cannot be $-1$ modulo
$M\ge16$. As $M>p+1$, the class of $p$ in $H_e$ therefore has
order exactly
\[
 \frac M{2^v}\ge\frac M{p+1}\ge\frac M{2x+1}
   >\frac43x^{\beta-1}>x^{\epsilon/4},\qquad
 \beta=\frac{64}{65}\left(2+\frac{799\epsilon}{100}\right).
                                                        \tag{2.7}
\]
Here $\beta-1>\epsilon/4$. At an unramified prime the source
Artin equality identifies this class with the image of its
$A$-coordinate, whose order is consequently at least as large.
This proves the long-orbit assertion of
Theorem~\ref{thm:empty-norms}. Its order is even, so the
$A$-coordinate has positive even order. The source predicate
\texttt{HEv} holds, and both its zero-coordinate and odd-order
\texttt{IsCs} alternatives are excluded.

\subsubsection*{The small-prime candidate and penalty conditions}

The source's lower transition scale is $p_0=k^2x^{26}$. Since
$k\ge5$ and $x>1$, $p_0>2x$. Its upper candidate scale is
$X'=\exp(x^{\epsilon/3})>2x$ for all sufficiently large $x$
with fixed $\epsilon>0$. Therefore, for an odd unramified prime
in \texttt{inertClass} and in the stated range, all the clauses
of \texttt{Cand} hold. In particular its conditional band clause
$p_0\le p\Rightarrow p\notin\mathrm{InBand}$ is vacuous.

The penalty condition also follows directly from the definitions.
Write $h_p$ for the $A$-coordinate. Because the prime is in
\texttt{inertClass}, the orbit step is $h_p$. The source chooses
one representative of each orbit using a map $b$ satisfying
$b(t)+j h_p=t$ for some $j$, and $b(t+h_p)=b(t)$.
These identities imply $b(b(t))=b(t)$ and that $b$ is constant
on each orbit. Distinct representatives therefore give disjoint
orbits. If $N_t$ counts top labels in the orbit of $t$, then
\[
               \sum_{t\in\mathrm{goodReps}}N_t\le |TT|.
                                                        \tag{2.8}
\]
Equivalently, the positions $0\le j<\operatorname{ord}(h_p)$
with $t+jh_p\in TT$ are in bijection with the top labels whose
representative is $t$. This verifies the equality of the two
top-count definitions used by the source, rather than assuming it.

Put $R_\epsilon=\lceil8(4+8\epsilon)/\epsilon\rceil+1$,
the source's constant \texttt{rG}, and let $c_{\mathrm{II}}$
be its fixed local constant \texttt{kII}. Below $p_0$ the
definition of the orbit penalty, independently of either word,
is
\[
 \kappa_t=\min(1,c_{\mathrm{II}}N_tR_\epsilon^2k x^{-26}),
 \qquad \Sigma=m\sum_{t\in\mathrm{goodReps}}\kappa_t.
                                                        \tag{2.9}
\]
Using (2.8), $m\le2n\le k$, and $|TT|\le n\le k$, gives
\begin{align}
 \Sigma&\le\max(0,c_{\mathrm{II}})R_\epsilon^2 m k|TT|x^{-26}
       \le\max(0,c_{\mathrm{II}})R_\epsilon^2 k^3x^{-26}\notag\\
 &\ll_{\epsilon,c_{\mathrm{II}}}
       x^{-25/2+(3003/100)\epsilon}\le x^{-\epsilon}
                                                        \tag{2.10}
\end{align}
for all sufficiently large $x$, by the upper \texttt{SizeK}
bound and $\epsilon<1/52$. No positivity of
$c_{\mathrm{II}}$ is needed for this upper bound.
The definition of \texttt{VV} now makes it exactly
\texttt{goodReps}, for every pair of word functions.

For completeness, the upper support scale is
$X''=xn\exp(x^{\epsilon/3})>2x$; here the lower degree window
gives the integer $n\ge2$. In the source definition
\texttt{NoGain}, the low-order alternative is excluded by
(2.7), the band alternative by $p<p_0$, the excessive-penalty
alternative by (2.10), the large-prime alternative by $p<X''$,
the prime-class alternative by \texttt{inertClass}, and the
odd-order alternative by \texttt{HEv}. Thus \texttt{NoGain}
is false. This proves the remaining assertions of the opening
theorem. It establishes neither a nonempty good-orbit set nor
a lower bound for the gain, and does not supply a cubic
character weighting.

The literal \texttt{FamCore} quantities \texttt{gainTot} and
\texttt{ioTot} sum local terms by \emph{rational residue
characteristic}. They do not select prime ideals of absolute norm
at most $2x$. The local data identify an orbit length with the
order of the $A$-coordinate of Frobenius, and the growth argument
explicitly uses orbits of order at least $x^{\epsilon/4}$.
These local quantities may therefore be nonzero when every ordinary
Hecke coefficient considered above vanishes. Theorem~\ref{thm:empty-norms}
rules out the ordinary low-norm comparison, including its signed
mixtures; it does not rule out that rational-prime orbit statistic.

At a fixed scale \texttt{FamCore} holds the frame, the Artin map,
the root discriminant, the auxiliary level primes, and the labelled
top set fixed across its family members. Varying the family index
does not create a different low-norm Euler support. Its covering
condition varies weight and translation data. A new comparison
must therefore identify the local orbit gain, with its
normalization, rather than substitute an ordinary Hecke sum for it.

These conclusions are ordinary mathematical consequences of the
displayed source definitions. No compilation or kernel verification
of the present implications in Lean is asserted.
